% !TEX root = ../main.tex

\section{Related Work}
\label{sec:related-work}

Treating predictive quality and the number of selected features as separate objectives makes both the quality of the returned subsets and the range of their sizes relevant to a comparison~\cite{Xue2016Survey,Jiao2024Survey}. The nondominated sorting genetic algorithm II (NSGA-II) provides the selection framework used in this study~\cite{Deb2002NSGAII}. It ranks candidates by dominance: one candidate dominates another if it is no worse in either objective and better in at least one. It selects the next population from parents and offspring together, using crowding distance to favor candidates in less crowded parts of each front. Crowding distance measures spacing in the objective values. These selection rules can be combined with different initializations, crossover operators, and mutation operators, so the behavior of a particular configuration also depends on those choices.

The number of selected features, also called subset cardinality, is an established consideration in population initialization. Previous work compares random and informed starts, allocates candidates across size ranges, and supplies sparse candidates, meaning subsets with few selected features~\cite{Deniz2019,Deniz2022,Kropp2022}. These choices affect which subset sizes are available before evolutionary search begins. One explicit approach first samples a subset size uniformly within a given range and then selects that many features without replacement~\cite{Miyandoab2024}. Miyandoab et al. compare NSGA-II with and without this initialization as well as their complete method, so a controlled comparison of initializations within one algorithm is itself prior art, and a later study extends the comparison to several population-based methods~\cite{Miyandoab2026}. Bounded starts are also used in biomarker discovery: Cong et al.~\cite{Cong2026} attribute their start with at most 50 features to Cattelani and Fortino~\cite{Cattelani2022}. The bounded initialization used in our comparison follows the size sampling of Miyandoab et al. with a smaller range. Its role is to provide access to small subsets under the chosen budget; the sampling procedure itself is existing work.

Opposition-based initialization seeks diversity by pairing candidates with their opposites~\cite{Tizhoosh2005OBL}. In the binary representation used by NSGAII-MIIP, the opposite selects every feature excluded by the original candidate and excludes every feature it selected~\cite{Li2025}. With independent random inclusion at probability one half, both candidates have an expected size of half the available features. Complementary pairs therefore do not allocate the population evenly across subset sizes or guarantee the presence of small subsets. This distinction matters when comparing opposition-based and bounded sparse starts.

Maintaining small subsets during search also depends on crossover, mutation, and the rule used to retain candidates. Methods that combine size-based initialization with adaptive crossover address both the initial population and the sizes of its descendants~\cite{Xue2022,Wang2026}. FASTENER, an evolutionary wrapper developed for land-cover classification from Earth observation data, uses feature information to guide crossover toward smaller subsets and keeps the search population separate from the stored Pareto front~\cite{Koprivec2020}. SparseEA and its feature-selection extensions combine initialization with feature scoring and search operators designed for sparse solutions~\cite{Tian2020SparseEA,Chu2024}. Selection can also preserve different subset sizes by retaining candidates from separate size groups~\cite{Li2026CORAL}. Together, these approaches show why an experiment that changes only the initial population differs from one that also changes how candidates are generated and retained.

The expected change in subset size and the amount of change produced by mutation are separate properties. Controlled mutation can balance the expected numbers of removed and added features~\cite{Oh2004,Li2025}, which keeps the expected subset size unchanged but leaves open how often any feature changes. Under the mutation rate examined here, that frequency depends on the current subset size relative to the total number of features.

An informed improvement phase provides another way to change the subsets produced by genetic operators. NSGAII-MIIP combines NSGA-II with an improvement phase guided by mutual information (MIIP)~\cite{Li2025}. Mutual information measures how much knowing one variable reduces uncertainty about another. The phase uses this information to guide additions, removals, and exchanges of features according to their relevance to the class and their redundancy with selected features. The resulting candidates are ranked together with the original candidates. Our comparison uses the genetic operators of this method with the two initializations and includes our implementation of the complete method with MIIP. This separates the change in initialization from the additional search performed by the improvement phase.

A finite budget makes access to small subsets an empirical question even when the initialization can generate every subset size, because sampling from a broad range does not guarantee that a finite population contains a candidate within a particular small-size limit. Existing work addresses this through size sampling and later insertion of new candidates~\cite{Miyandoab2024} and compares sparse and dense starts under a fixed evaluation target~\cite{Cong2026}, so compact-front access is already a subject of experimental study.

The evaluation measure determines how differences in subset-size coverage appear in the reported result. FASTENER introduced the area under the front (AUF), which averages the best predictive score available with at most a given number of features over a range of feature-count limits~\cite{Koprivec2020}. A subset with three features can therefore contribute at limits of three, four, or more features; a missing subset of exactly four features does not by itself imply a missing score at that limit. Hypervolume measures the area between the front and a chosen reference point in this two-objective setting. Its value depends on both the extent of the front and the objective values, as well as the reference point~\cite{Auger2012}. We use the AUF definition from FASTENER.

Repeated candidates connect search behavior to computational cost. Duplicate handling has been used to address lost diversity in evolutionary feature selection~\cite{Jiao2024PRDH,Wang2023Dup}, while caching stores evaluation results so that the same subset need not be fitted again~\cite{Li2020GADMS}. Removing duplicates changes the candidates considered by the search; reusing a stored result saves a fit. A run can consequently generate many candidates while consuming few new model fits. Reporting proposals, new fits, and elapsed time together follows the broader distinction between evaluation budgets and computational time in optimization benchmarking~\cite{BartzBeielstein2020}. It also helps distinguish productive search from repeated requests for previously evaluated subsets.

These studies establish the importance of initialization, subset-size control, front evaluation, and repeated candidates; our focus is how these factors affect the interpretation of a comparison within a specified configuration.

\subsection{Our Contributions in Relation to Existing Work}
\label{sec:positioning}

Table~\ref{tab:positioning} lists the contributions of this paper in order of importance. For each, it states what is new here and which prior work it builds on.

\begin{table}[!tbp]
\caption{Contributions of this paper, in order of importance, and the prior work each builds on}
\label{tab:positioning}
\centering
\small
\setlength{\tabcolsep}{4pt}
\begin{tabular}{@{}>{\raggedright\arraybackslash}p{22mm}>{\raggedright\arraybackslash}p{52mm}>{\raggedright\arraybackslash}p{44mm}@{}}
\toprule
Contribution & New in this paper & Builds on \\
\midrule
Reach, not quality, explains the AUF gap
  & On 46 datasets, 39 of them attempted with three classifiers, reach accounts for 98--100\% of the AUF difference between the OBL and the sparse start, mostly at limits that the OBL start never reaches. At sizes both reach, the group mean of test F1 differs by less than 0.02 under two classifiers, and under the third only without the smallest test sets.
  & Initialization-only comparison within NSGA-II~\cite{Miyandoab2024}; sparse and dense starts compared for access to compact fronts in one gene-panel application~\cite{Cong2026}; AUF~\cite{Koprivec2020}, used unchanged. \\
\addlinespace
When reach fails
  & The probability that an initial population contains a small subset, and an observed boundary: within the benchmark budgets, the opposition-based search reaches ten features on every dataset with at most 99 features and on none with more than 512. With ten times the budget, it reaches them on all four pilot datasets.
  & Opposition-based initialization~\cite{Tizhoosh2005OBL,Li2025}; benefits of sparse sampling in large sparse problems~\cite{Kropp2022}. \\
\addlinespace
An informed improvement phase does not close the gap
  & The mutual-information phase, in our implementation of NSGAII-MIIP, raises the share of runs that reach ten features by only 0.04--0.08.
  & NSGAII-MIIP~\cite{Li2025}. \\
\addlinespace
The cost of a sparse start
  & Neither start covers all subset sizes. The sparse start's best test F1 at any size is lower by 0.007--0.029 on average, and its fronts end by 36 features. Starts over all sizes do not recover the large subsets; under the tested operators they disappear during the search.
  & Size sampling over a chosen range~\cite{Miyandoab2024}, used here with 1--30 features; size-allocated and bounded starts~\cite{Deniz2019,Cattelani2022}. \\
\addlinespace
Why the search stays near its start
  & The expected number of changed features under balanced mutation, $2k/d$, with unchanged children at approximately the rate $e^{-2k/d}$, and the observation that 96--99\% of the subsets proposed in sparse runs had already been evaluated. A forced change lowers their cost, the repeats only partly, and not the quality.
  & Balanced mutation~\cite{Oh2004} and its unchanged expected size~\cite{Li2025}; duplicate handling and caching~\cite{Jiao2024PRDH,Wang2023Dup,Li2020GADMS}. \\
\addlinespace
Reporting recommendations
  & Quantities to report with AUF at a feature limit: reach, scores at sizes both methods reach, the best score at any size, proposals, new fits, time, and unfinished runs.
  & Evaluation budgets versus computing time in benchmarking~\cite{BartzBeielstein2020}. \\
\bottomrule
\end{tabular}
\end{table}

Two studies are closest to ours. Miyandoab et al.~\cite{Miyandoab2024} described the size sampling that our sparse start uses, over all sizes and with later insertion of new candidates, and compared it with the default start within NSGA-II by hypervolume. We use it over 1 to 30 features, without insertion, as an experimental condition and not as a method of our own. Cong et al.~\cite{Cong2026} compare sparse and dense starts at a fixed evaluation target and report access to compact fronts in one gene-panel application. What we add is not the comparison of starts itself but its reading: the split of an AUF difference into reach and quality at matched sizes, the dependence of reach on the number of features and the budget, the comparison with an informed improvement phase, and the counts of proposals and fits that explain the cost, across 46 datasets and up to three classifiers. The initialization, the mutation, and AUF are existing methods.

We test NSGA-II with the genetic operators of Li et al.~\cite{Li2025}, in our implementation, with and without their mutual-information improvement phase, so the results describe these configurations and the tested budgets and classifiers, not NSGA-II in general. Methods that change crossover or survival to keep small subsets~\cite{Xue2022,Wang2026,Tian2020SparseEA,Chu2024,Li2026CORAL} address the same problem by other means, and we do not compare against them. The returned front shows which subset sizes survive to the end of a run. It does not show which sizes were visited during the search, which would require a record of all evaluated subsets.
